\section{Experimental Setup}
\label{sec:setup}

\subsection{Dataset}

We use the public training portion of the PhysioNet/CinC 2016
corpus~\cite{liu2016open,goldberger2000physiobank}: 3240 mono recordings,
\SI{20.2}{\hour} in total, sampled at \SI{2}{\kilo\hertz}, with durations from
about \SI{5}{\second} to \SI{120}{\second}. Labels are per recording, and
\SI{20.5}{\percent} of recordings are abnormal, an imbalance ratio of
\num{3.87}:1. The corpus is divided into six sub-databases, \texttt{training-a}
through \texttt{training-f}, each collected by a different group at a different
site with different equipment. Abnormal prevalence differs across them by
\SI{68.9}{pp}, which is the single most important property of this corpus for
the present study.

All results below therefore come from a self-defined split of the six
subdataset, and published challenge scores are context rather than a
like-for-like comparison. We restate this wherever those numbers appear.

\subsection{Splits}

Splits are made at the \emph{recording} level, never at the window level, and
stratified jointly on label and sub-database. Recording-level splitting prevents
windows of the same patient --- sharing a stethoscope, a noise floor and a
murmur --- from appearing in more than one partition; this is the most
consequential guard in the whole protocol. Joint stratification keeps both the
class balance and the site mix constant, so that the test partition cannot end
up dominated by one clinic. The realised deviation is at most \SI{0.07}{pp} in
class balance and \SI{0.24}{pp} in site composition.

Windowing yields \num{30695} training, \num{3329} validation and \num{3362} test
windows, the latter drawn from 484 recordings of which \SI{20.5}{\percent} are
abnormal. No recording was excluded for quality or loading failure.

\subsection{Parameters}

Features use a 256-point Hann window with hop 32 (\SI{16}{\milli\second}), 32
mel bands over 25--\SI{400}{\hertz}. MFCC uses 13 coefficients with $\Delta$ and
$\Delta^2$ (width 5) and six aggregation statistics, giving 234 dimensions;
log-Mel windows are $32\times188$; PWP uses a level-6 db4 tree grouped into 12
perceptual bands with 7 descriptors each, giving 84 dimensions. All classical
features are standardised with statistics fitted on the training split only. No
degenerate all-zero mel band and no near-constant feature was produced.

Classical hyperparameters were chosen by five-fold stratified group cross
validation on the training split, grouping by recording and scoring balanced
accuracy. The selected SVM uses $C = 10$, $\gamma = \num{5e-4}$, with
cross-validated MAcc $0.874 \pm 0.008$. The CNN was trained for at most 60
epochs with Adam at $10^{-3}$, batch size 64, dropout \num{0.3}, plateau
scheduling and early stopping on validation MAcc; it stopped after 20 epochs
with its best epoch at 8 and validation MAcc \num{0.915}. Model selection,
hyperparameter search and decision thresholds used training and validation data
only; the test partition was read once, after selection was frozen.

\subsection{Metrics and statistics}

The primary metric is the official challenge metric, modified accuracy
$\MAcc = (\mathrm{Se} + \mathrm{Sp})/2$, which is identical to balanced
accuracy. Plain accuracy is reported but is not an appropriate headline: with
roughly \SI{80}{\percent} normal recordings, a model that always predicts
\emph{normal} scores \num{0.795} accuracy at \num{0.500} MAcc. Trivial baselines
appear in Table~\ref{tab:baselines} so that the reader need not do this
arithmetic. Sensitivity, specificity, precision, F1, ROC-AUC and the Matthews
correlation coefficient are also reported.

Confidence intervals are obtained by bootstrap over \emph{recordings}, which is
the correct unit of independence; resampling windows would treat the several
windows of one recording as independent observations and would produce intervals
far too narrow. Pairwise model comparison uses McNemar's
test~\cite{dietterich1998approximate}, because all models are evaluated on the
same recordings and the samples are therefore paired, with Holm--Bonferroni
correction across the family of ten comparisons.
